\chapter{Résultats}

\section{Types de résultats}

Après analyse, le logiciel peut afficher :

\begin{itemize}
  \item déplacements nodaux ;
  \item réactions d'appui ;
  \item efforts internes dans les barres ;
  \item diagrammes d'efforts ;
  \item déformée ;
  \item résultats plaques lorsque disponibles ;
  \item enveloppes de résultats selon les fonctionnalités actives.
\end{itemize}

\section{Tableaux}

Les tableaux de résultats sont accessibles depuis le menu \menu{Résultats}. Ils
permettent de consulter les valeurs numériques par cas de charge ou combinaison.

\section{Diagrammes de barres}

Les diagrammes disponibles sont :

\begin{itemize}
  \item effort normal N ;
  \item efforts tranchants Vy et Vz ;
  \item moments My et Mz ;
  \item torsion T.
\end{itemize}

\section{Déformée}

La déformée permet de visualiser qualitativement le comportement global du
modèle. Elle doit être lue comme une aide graphique ; les valeurs numériques
restent dans les tableaux.

\section{Résultats plaques}

Lorsque les surfaces sont prises en charge, les résultats plaques peuvent être
consultés sous forme de tableaux ou de cartes de contours.

\section{Conventions de signe}

Les conventions de signe dépendent du repère local des éléments et du moteur de
calcul. Le logiciel centralise progressivement ces conventions afin de rendre la
lecture plus stable.

\todoDoc{Ajouter un schéma des axes locaux et des signes positifs pour N, V, M et T.}
